\section{Plan de trabajo}
\setlength{\columnseprule}{0pt}
\onecolumn
\vspace{-0.3cm}
\begin{flushleft}
\resizebox{0.9\textwidth}{!}{%
\begin{ganttchart}[
    x unit=0.9cm, y unit title=0.5cm, y unit chart=0.65cm,
    vgrid={*1{dotted}}, hgrid,
    title height=1, title label font=\bfseries\small,
    bar/.style={fill=colorH},
    bar incomplete/.style={fill=colorH!25},
    group/.append style={fill=gray!50, draw=gray!50},
    progress label text={},
    bar height=0.7,
    bar top shift=0.15,
    group right shift=0.15, group top shift=0.15, group height=.7,
    bar label font=\small,
    group label font=\small,
    bar label node/.append style={anchor=east, align=left, text width=12cm},
    group label node/.append style={anchor=east, align=left, text width=12cm},
    bar inline label node/.append style={font=\small\bfseries\color{white}}]{1}{6}
  \gantttitle{May}{1}\gantttitle{Jun}{1}\gantttitle{Jul}{1}
  \gantttitle{Ago}{1}\gantttitle{Sep}{1}\gantttitle{Oct}{1}\\
  \ganttgroup[inline=false]{\textbf{Protocolo de consenso con PVA}}{1}{2}\\
  \ganttbar[progress=80,inline=false]{Implementar protocolo de consenso con seguimiento de líder}{1}{2}
  \ganttbar[progress=80,inline=true]{\textbf{80\%}}{1}{2}\\
  \ganttbar[progress=50,inline=false]{Integrar PVA al protocolo de consenso}{1}{2}
  \ganttbar[progress=50,inline=true]{\textbf{50\%}}{1}{2}\\
  \ganttbar[progress=50,inline=false]{Validar en simulación estática y dinámica}{2}{2}
  \ganttbar[progress=50,inline=true]{\textbf{50\%}}{2}{2}\\
  \ganttgroup[inline=false]{\textbf{Comparativa y métricas}}{3}{4}\\
  \ganttbar[progress=0,inline=false]{Implementar campos de potencial como baseline}{3}{3}
  \ganttbar[progress=0,inline=true]{\textbf{0\%}}{3}{3}\\
  \ganttbar[progress=0,inline=false]{Definir y calcular métricas de desempeño}{3}{4}
  \ganttbar[progress=0,inline=true]{\textbf{0\%}}{3}{4}\\
  \ganttbar[progress=0,inline=false]{Comparar PVA vs campos de potencial}{3}{4}
  \ganttbar[progress=0,inline=true]{\textbf{0\%}}{3}{4}\\
  \ganttgroup[inline=false]{\textbf{Redacción de tesis y experimentación en físico}}{4}{5}\\
  \ganttbar[progress=0,inline=false]{Redactar tesis incorporando resultados finales}{4}{5}
  \ganttbar[progress=0,inline=true]{\textbf{0\%}}{4}{5}\\
  \ganttgroup[inline=false]{\textbf{Cierre y defensa}}{6}{6}\\
  \ganttbar[progress=0,inline=false]{Integrar resultados experimentales a la tesis}{6}{6}
  \ganttbar[progress=0,inline=true]{\textbf{0\%}}{6}{6}\\
  \ganttbar[progress=0,inline=false]{Conclusiones y defensa}{6}{6}
  \ganttbar[progress=0,inline=true]{\textbf{0\%}}{6}{6}\\
\end{ganttchart}
}
\end{flushleft}
